\chapter{Deep Dive into Data-Sources: Sensors and how they work}\label{title:data-sources}
Sensors are fundamental data sources in many technical systems, as they allow devices to observe and measure physical phenomena. By converting environmental or mechanical changes into electrical signals, sensors provide the raw data required for analysis, monitoring and control.

This chapter presents several important sensor types and explains the principle behind their operation. It also introduces how sensor signals are transformed into digital values and which factors influence the quality of the resulting measurements.

\section{IMUs}\label{title:imu}

An IMU\footnote{Inertial Measurement Unit} is a device that measures the angular acceleration and linear acceleration. To achieve this, an IMU incorporates an accelerometer~\ref{title:accelerometer}, a gyroscope~\ref{title:gyroscope} and optionally a magnetometer. If an IMU uses a magnetometer, it is referred to as IMMU\footnote{Inertial Measurement \& Magnetometer Unit}

\subsection{Accelerometer}\label{title:accelerometer}
An accelerometer is a sensor that measures the proper acceleration of a physical object. The proper acceleration being the acceleration of an object relative to an observer who is in free fall.

When looking at a resting accelerometer, it is important to note that it will usually fluctuate around $-9.81$. This is due to gravity acting on the sensor.

An accelerometer distinguishes between \textbf{static} and \textbf{dynamic} acceleration. Static acceleration always acts on the sensor (e.g.\ gravity or friction), whilst dynamic acceleration is non-uniform (e.g.\ vibration or shock).\\ 
\noindent\cite{LRWeb:29}

\boldtitle{How does an accelerometer work?}

An typical MEMS accelerometer contains four springs, multiple fixed plates (electrodes), and an H-shaped proof mass that is tied to the springs (allowing it to move).

The main way an accelerometer works is through the concept of capacitance. Capacitance is defined as:

\[
C = \frac{\varepsilon A}{d}
\]

$C$ = Capacitance (Farads)\\
$\epsilon$ = Dielectric permittivity\\
$A$ = Plate overlap area (square meter)\\
$d$ = Distance between plates (meter)\\

To understand how an accelerometer works, we focus on the distance. The greater the distance, the smaller the capacitance, and vice versa. When we now take a look at a differential capacitor, which is made up of two stationary plates ($P_1$ and $P_2$) and a movable proof mass ($M$).

When moving $M$ towards $P_1$, the capacitance between $M$ and $P_1$ gets larger, while the capacitance towards $M$ and $P_2$ gets smaller and vice versa. This now lets us calculate the difference between those values, which can then be processed by the sensor's electronics to obtain the measured acceleration.

To measure acceleration for the other two axis, we have to mirror the differential capacitor twice, so that they are at a 90° angle towards each other.

\noindent\cite{LRWeb:28}

\begin{figure}[H]
    \centering
    \includegraphics[width=0.5\linewidth]{images/Accelerometer.jpg}
    \caption{MEMS Accelerometer}\label{fig:accelerometer}
\end{figure}

\noindent\cite{LRWeb:30}


\subsection{Gyroscope}\label{title:gyroscope}

A gyroscope is a device that measures angular velocity or maintains orientation.\\
\noindent\cite{LRWeb:31}

A typical MEMS gyroscope takes advantage of the Coriolis effect.

\boldtitle{Coriolis effect}

The Coriolis effect is a phenomenon that occurs on a moving body on which an angular rate is applied. This will apply a force on the body depending on the angular velocity, causing a perpendicular displacement of the mass.

\begin{figure}[H]
    \centering
    \includegraphics[width=0.5\linewidth]{images/coriolis.png}
    \caption{Coriolis effect}\label{fig:coriolis}
\end{figure}

\noindent\cite{LRWeb:30}

\boldtitle{How does a gyroscope work?}

The way a typical MEMS gyroscope works is similar to how a typical MEMS accelerometer works. The difference is that the mass is constantly moving if angular velocity is now applied, the mass creates a different capacitance based on the rotation.

\begin{figure}[H]
    \centering
    \includegraphics[width=0.5\linewidth]{images/gyroscope.jpg}
    \caption{Visualization of a gyroscope}\label{fig:gyroscope}
\end{figure}

\noindent\cite{LRWeb:30}

\subsection{Sensor Fusion (IMU-Integration)}
Sensor fusion is the algorithmic combination of measurements from an Inertial Measurement Unit (IMU) and, optionally, additional sensors (e.g., GPS, magnetometer, vision), in order to estimate system states such as orientation, velocity, and position with higher accuracy and robustness than any individual sensor can provide.

Sensor fusion is usually used to:
\begin{itemize}
    \item Compensate for sensor bias and integration drift~\ref{title:drift}
    \item Reduce measurement noise through probabilistic filtering~\ref{title:noise-removal}
    \item Improve state estimation consistency and reliability
    \item Exploit complementary characteristics of different sensors
\end{itemize}

\noindent\cite{LRWeb:32}

There are a lot of IMU Sensor Fusion algorithms, the most important ones are:
\begin{itemize}
    \item Any filter from the Kalman Filter Family~\ref{title:kalman}
    \item Madgwick / Mahony Filter~\ref{title:madgwick}
    \item Complimentary Filter
    \item Machine Learning approach
\end{itemize}

\noindent\cite{LRWeb:32}

\section{Sensors for collecting environmental data}
Many applications require sensors that monitor environmental conditions such as air quality or humidity. These measurements are important for areas such as environmental monitoring, smart buildings, and health-related systems.

The following subsections present common sensors used to collect environmental data and explain the principles behind their operation.

\subsection{Fine dust-Sensors}
Fine particulate matter (PM) is a critical air quality parameter due to its documented effects on human health. In low-cost air quality monitoring systems, optical laser-based sensors are commonly used to estimate particulate matter concentrations. Among these, the \textit{SDS011} and the \textit{PMS7003} are two widely deployed sensors. Both devices operate on the principle of laser light scattering, yet differ in internal design, data processing, and output characteristics.

\boldtitle{SDS011}
The SDS011 detects fine-dust matter through laser light scattering. The particles are sucked out of the environment through an integrated fan. In the interior of the fan, there is a laser which, in cooperation with a photo-diode, counts the amount of particles in the air. The sensors outputs two different values, these being the concentration of PM2.5 (Particles smaller than 2.5µm) and PM10 (Particles smaller than 10µm).

\noindent\cite{LRWeb:35}

\boldtitle{PMS7003}
The PMS7003 also detects fine-dust matter through laser light scattering. The sensor uses a small laser to shine light on the dust particles in the air. When the laser hits a particle, the light scatters. A detector then measures the scattered light, and the sensor's processor estimates the amount of particles and their size. The calculated concentrations are PM2.5 and PM10.

\noindent\cite{LRWeb:34}

\subsection{Humidity-Sensors}\label{title:humidity}
Humidity is an important environmental parameter that affects our lives in many ways. High humidity can negatively affect our health by impairing the body's ability to cool down, whilst low humidity can dry out sensitive skin. Furthermore, high humidity can cause mold to form, thus it's important to have a view on the current humidity. 

\noindent\cite{LRWeb:37}

There are multiple types of humidity sensors among them are the \textit{capacitive} the \textit{resilient} and the \textit{thermally conductive humidity sensor}.

\boldtitle{Capacitive humidity sensor}
Capacitive humidity sensors are the most commonly used sensors of their type. They consist of two electrodes with a special material in between called a dielectric.

This material absorbs moisture from the air. When the air becomes more humid, the material absorbs more water vapor, causing the dielectric constant to increase, which then increases the capacitance of the sensor.

\boldtitle{Resilient humidity sensor}
Resilient humidity sensors measure humidity by detecting changes in electrical resistance. This sensors incorporates a hygroscopic, conductive layer with interlocking electrodes.

When the air becomes more humid, the material absorbs water vapor. This increases its electrical conductivity, which means the resistance decreases. By measuring this change in resistance, it's now possible to calculate the relative humidity.

\boldtitle{Thermally conductive humidity sensor}
Thermally conductive humidity sensors measure absolute humidity. They use two probes, where one is exposed to the surrounding air, and the other is sealed with dry nitrogen.

The sensor then compares the resistance difference between the two probes. This difference is directly related to the amount of water vapor in the air, allowing the sensor to determine the absolute humidity.

\noindent\cite{LRWeb:36}

\section{Measurement Methods: From the physical unit to the digital value}
Before sensor data can be processed by digital systems, the measured physical quantities must be converted into digital values. This process involves sampling, quantization, and evaluation of measurement quality.

The following subsections introduce the basic principles of analog-to-digital conversion and important sensor characteristics such as resolution, repeatability, and accuracy.

\subsection{Analog to Digital Transformation}
ADCs\footnote{Analog-to-Digital Converters} are a component when dealing with the continuous (analog) signal emitted from a sensor. The purpose of the ADC is to break down the continuous signal into a sequence of discrete values.
\begin{figure}
    \centering
    \includegraphics[width=0.75\linewidth]{images/ADC.png}[H]
    \caption{An analog signal transformed into a digital signal}\label{fig:adc}
\end{figure}

The way an ADC works is through the sampling procedure. For this, the parameter sampling rate is necessary. The sampling rate determines in which interval the data are obtained from the analog stream. This parameter may be expressed as samples per second (S/s) or in hertz (Hz).

When choosing a sampling rate, there are a few general rules of thumb. The first one being that it should not be too low, as the digital values are not accurate to the real analog values anymore. As can be seen in figure~\ref{fig:adc-low}. The second thing to consider is that a too high sampling rate can lead to unnecessarily high power usage.To find a suitable sampling rate, the Nyquist theorem can be followed. This theorem says that the sampling rate needs to be twice as large as the highest frequency in the analog signal.

\begin{figure}
    \centering
    \includegraphics[width=0.5\linewidth]{images/ADC-Low.png}
    \caption{Analog signal compared to digital signal with a low sampling rate}\label{fig:adc-low}
\end{figure}

\noindent\cite{LRWeb:41}

\subsection{Sensor Resolution, Repeatability \& Accuracy}
\boldtitle{Resolution}
Sensor resolution is described as the smallest possible change that a sensor can reliably detect and represent in its output. Higher resolution enables finer discrimination of changes in the input signal and therefore provides greater measurement detail. However, increased resolution typically comes with higher cost and does not automatically improve measurement accuracy or usefulness. In many applications, the downstream processing or control system cannot exploit the additional detail, regardless of the ADC, making the higher resolution economically unjustified.

\noindent\cite{LRWeb:51, LRWeb:52}

\boldtitle{Repeatability}
The repeatability of the sensor represents the consistency of the data obtained from the same sensor operating under similar conditions. In more technical words:

\begin{quote}
    [\dots] repeatability (or sometimes referred to non-repeatability), is the maximum difference between transducer output readings for repeated loadings under identical loading and environment conditions.
    Repeatability is typically expressed as a standard deviation or a percentage of the full-scale measurement range. 
\end{quote}

\noindent\cite{LRWeb:52}

Conditions that negatively affect repeatability include:
\begin{itemize}
    \item Mechanical tolerances
    \item Temperature variations
    \item Electronic noise
    \item Calibration errors
\end{itemize}

\noindent\cite{LRWeb:52}

\boldtitle{Accuracy}
Sensor accuracy is defined as `how close the measured value is to the true value'. Technically, the accuracy is described as the average deviation between the real output versus the theoretical output. Accuracy is often treated as a key indicator of a sensor's quality, but in reality it is a composition of accuracy, repeatability and resolution. 

Accuracy is affected by:
\begin{itemize}
    \item Offset (bias)
    \item Gain Error
    \item Nonlinearity
\end{itemize}

Environmental factors like temperature variation, aging and mechanical stress can further degrade accuracy over time. Especially for accuracy, a good calibration play a key role in maximizing the quality of this factor.

\noindent\cite{LRWeb:52}